\chapter{Use of AI}\label{sec:use_of_ai}

During this thesis project, I used AI-based assistants, including OpenAI GPT-5.3 Codex\footnote{OpenAI, ``GPT-5.3 Codex,'' \url{https://openai.com}.} and Claude Sonnet 4.6\footnote{Anthropic, ``Claude Sonnet,'' \url{https://www.anthropic.com/claude}.}, as support tools in software development. These tools were mainly used to speed up drafting and refinement of selected code segments related to the thesis implementation, for example by proposing initial structures, suggesting refactorings, and helping identify potential errors. AI-generated code was never integrated blindly: each suggestion was reviewed in context, tested, and revised where necessary to satisfy project requirements and maintainability goals.

I also used the same AI tools to support parts of the writing process for this document, such as improving phrasing, reorganizing draft text, generating preliminary adjustments for specific sections, and translating the Abstract section into German. All generated text was critically checked against the intended meaning, technical content, and academic standards. Final wording, structure, and claims were manually edited by me to ensure consistency, correctness, and alignment with the overall thesis. In addition, the ``Principles for the Use of Artificial Intelligence (AI) by Students at the University of Stuttgart'' were followed throughout the work.\footnote{University of Stuttgart, ``Principles for the Use of Artificial Intelligence (AI) by Students at the University of Stuttgart,'' \url{https://www.student.uni-stuttgart.de/digital-services/dokumente/ki-principles-for-students-en.pdf}.}
